\section{Outlook}\label{sec:outlook}

\subsection{What is settled}

Three things are now closed, and it is worth separating them from what is not.

\textbf{The generator and the transfer are identified.}
Theorem~\ref{thm:kernel} writes the form in closed shape and shows the whole
shift dependence is one hyperbolic factor; Proposition~\ref{prop:unimodular}
shows the transfer is an all-pass filter, exactly and unconditionally. Everything
in Sections~\ref{sec:allpass}--\ref{sec:region} follows from the second.

\textbf{Two of the program's own facts are explained rather than recorded.} The
density identity $b(\tau^2)+w_0=2\theta'(\tau)$ is the statement that the Weil
symbol is a group delay (Proposition~\ref{prop:groupdelay}), and the band edge at
$\gamma_1$ of the seventh selection rule is the statement that the cheapest way
to avoid every zero is to live below the lowest one (Section~\ref{sec:region}).

\textbf{The shift--length plane is settled.} Corollary~\ref{cor:graded} makes the
shifted family a graded criterion, one zero-free strip per shift; the contraction
region has no boundary under the Riemann hypothesis, so the critical path sought
by \cite{CriticalPath} does not exist; and Remark~\ref{rem:disproof} says what
the plane is for instead.

\textbf{The delay test is run, and it is archimedean.} No correlator of a
unitary defect theory can wind at all (Proposition~\ref{prop:nowinding}); the
protected sector of the Wilson-line defect CFT gives exactly zero
(Corollary~\ref{cor:protected}); and the clause is satisfied by the free
dilatation generator on a half-line, up to the constant $\log\pi$. What
discriminates is the expansion, whose leading coefficient is the degree
(Corollary~\ref{cor:degree}) and whose $\tau^{-2}$ coefficient forces the
archimedean shift to within $1/2\sqrt3$ of one half
(Corollary~\ref{cor:signrule}) --- excluding Liouville at every coupling.

\textbf{The margin is identified, is positive, and is not arithmetic.} It is the
lower sampling constant of the zeros on $PW_{L/2}$, with no upper one
(Proposition~\ref{prop:noupper}); it is strictly positive because the group delay
diverges (Proposition~\ref{prop:positive}); and its superexponential collapse is
reproduced, to within a ratio of logarithms tending to one, by an arithmetic-free
point set of the same counting function (\S\ref{subsec:density}). That closes the
question Section~\ref{sec:region} left open, in the sense of answering it rather
than of resolving it in the direction expected.

\subsection{What is open, ranked}

\begin{enumerate}
\item \textbf{Find an amplitude with the shift at a quarter.}
Corollary~\ref{cor:signrule} is the sharpest instrument this paper has, and it is
cheap: a candidate's reflection amplitude must contain a $\Gamma$ whose shift lies
within $1/2\sqrt3$ of one half. Every amplitude in the reflection literature is
built from $\Gamma(\text{integer}+i\,\cdot)$ and fails. Liouville fails because
its reflection is off a wall on a half-line, which puts the shift at an integer;
an amplitude with the shift at a quarter would have to come from something whose
natural variable is $s/2$ --- a square root of the radial coordinate, or a
two-sheeted cover. That is a concrete structural hint and it is the first one this
program has had about the archimedean factor. Until it is followed, no candidate
can pass Condition~\ref{cond:bilinear}.

\item \textbf{What does the arithmetic control, if not the margin?}
\S\ref{subsec:density} measures the gap between the margin of the true zero set
and the margins of matched-density arithmetic-free sets: about three orders of
magnitude at $L=2$, growing like $\log L$ against a total growing like $e^L$. That
gap is a well-defined and computable quantity which \emph{does} carry arithmetic,
and neither this program nor the literature has studied it. It replaces the
previous second item, which asked how $\varepsilon_L$ collapses; that question now
has an answer with no arithmetic in it.

\item \textbf{The universal frequency profile of \S\ref{subsec:profile}.} The
obstruction's spectral content is a fixed profile in $\tau/\tau_c$, stable to about
one percent while the margin moves through $10^{80}$. That is the one structural
regularity the measurement turned up that is not explained here, and an exact
statement of it would be worth having.

\item \textbf{Problems~\ref{prob:trace} and \ref{prob:curvature}}, which
Section~\ref{sec:endpoint} shows are the two forms of one question.

\item \textbf{Does the vanishing theorem generalize?} Whether "BPS implies
vanishing expectation unless superconformal" extends beyond the class of
\cite{OWL}. If it does, it is a general obstruction of the same kind as this
program's own exclusions and belongs in the selection index; if it does not, the
exception is where to look.

\item Only after these, a specific theory. Section~\ref{sec:conditions} says what
it must satisfy, and Remark~\ref{rem:angle} says where to look first.
\end{enumerate}

\subsection{Scope}

Nothing above is proposed as a route to the Riemann hypothesis, and
Section~\ref{sec:region} explains structurally why no argument of this shape could
be one.

Written proofs: Proposition~\ref{prop:noupper}, Proposition~\ref{prop:positive},
Proposition~\ref{prop:defectdelay}, Proposition~\ref{prop:nowinding},
Proposition~\ref{prop:expansion}, Proposition~\ref{prop:rigidity},
Theorem~\ref{thm:kernel}, Proposition~\ref{prop:unimodular},
Proposition~\ref{prop:groupdelay}, Proposition~\ref{prop:resonance},
Proposition~\ref{prop:flux}, Lemma~\ref{lem:multiplicative},
Theorem~\ref{thm:abelian}, Proposition~\ref{prop:hyperbolic},
Proposition~\ref{prop:monotone}, Proposition~\ref{prop:dichotomy},
Corollary~\ref{cor:graded}, Proposition~\ref{prop:cayley},
Proposition~\ref{prop:endpoint}, Proposition~\ref{prop:algebra} and
Proposition~\ref{prop:flat}.

Conditional on the Riemann hypothesis: Proposition~\ref{prop:flux},
Proposition~\ref{prop:positive}, the reading \eqref{eq:sampling} of the margin as
a sampling constant, and the second case of Proposition~\ref{prop:dichotomy},
which uses the weaker hypothesis stated there. The computations of
Section~\ref{sec:sampling} are themselves unconditional: they evaluate the Weil
functional, which is defined without the hypothesis. Proposition~\ref{prop:resonance}'s passage from one zero to the sum is
stated rather than proved uniformly and is supported by
Appendix~\ref{app:checks}.

Proof sketch: Proposition~\ref{prop:trace}, resting on a smoothing estimate
quoted from \cite{Background} and on the Sobolev trace threshold.

Labelled numerical computation: every table of Section~\ref{sec:sampling}, each
entry a trial-space Rayleigh quotient and hence a one-sided bound, produced by
unregistered programmes recorded in the investigation's notes; and the
convergence rule of \S\ref{subsec:profile}, which is itself a measurement.

Heuristic, and labelled where it occurs: the single-parameter relation
$\kappa_L\sqrt{m_L}\approx\sqrt{Q[Xe_L]}$ of Proposition~\ref{prop:kappa}, whose
support is three data points --- independently reproduced to a few percent, per
Remark~\ref{rem:samplingstatus}, but still three; the detection scale $\omega\sim1/L$ of
Remark~\ref{rem:artifact}; and Remark~\ref{rem:principal}, which is a reading
resting on an identification the program declines to make.

Not claimed anywhere: a source, a positivity statement, a matched gauge theory,
an excluded gauge theory, or novelty for the idea that the zeros behave like
resonances.

\subsection*{Preparation}

This working manuscript was produced with a large language model, Anthropic's
Claude (model \texttt{claude-opus-5}), working directly in the investigation
repository: the derivations and proofs of
Sections~\ref{sec:generator}--\ref{sec:endpoint}, the conditions of
Section~\ref{sec:conditions}, the identification and measurement of
Section~\ref{sec:sampling}, the delay analysis of Section~\ref{sec:delay}, the
three check programs of Appendix~\ref{app:checks}, and the drafting throughout. The proposal the paper analyses is Edward Baker's, as is
the open Wilson line construction of \cite{OWL} and the earlier shifted-zeta
material of \cite{CriticalPath,CumulativePath} on which Section~\ref{sec:region}
builds; Remark~\ref{rem:attribution} records the priority of the hyperbolic
identity. The research notes accompanying the investigation name the model used
for each and record the exploratory computations deliberately outside the
registered checks. Every statement is classified in the text as a written proof,
a proof sketch, a labelled numerical computation, a heuristic, or a reading.
